\chapter{Getting Started}
\label{ch:getting-started}

This document is both an example thesis and the manual for the template that produced it. Everything you see here (the title page, the supervisory committee page, the table of contents, the page numbering) is generated by the \texttt{uvicphysastro} document class, following the formatting requirements of the University of Victoria's Faculty of Graduate Studies. The chapters themselves explain how to use the template, and give examples of content that astronomy theses commonly need: bibliographies in journal styles (Chapter~\ref{ch:bibliographies}), complex tables (Chapter~\ref{ch:tables}), and astronomical notation (Chapter~\ref{ch:astronomy}).
\section{UVic Formatting Requirements}
\label{sec:uvic-requirements}

UVic does not provide an official \LaTeX\ template. Instead, the Faculty of Graduate Studies publishes its requirements in two places:
\begin{itemize}
    \item Scope, structure and formatting: \url{https://www.uvic.ca/students/graduate/thesis-dissertation/scope-structure-and-formatting/index.php}
    \item Thesis format checklist and sample pages: \url{https://www.uvic.ca/graduatestudies/forms-policies/data/sample-samplepages.pdf}
\end{itemize}
If either link no longer works, start from the main thesis and dissertation page of the Faculty of Graduate Studies, \url{https://www.uvic.ca/students/graduate/thesis-dissertation/index.php#ipn-thesis-dissertation}, which links to the current versions.
The template follows the checklist and sample pages as last updated on 7 October 2025. Responsibility for meeting the formatting standards still rests with you, so read through both before you submit. Chapter~\ref{ch:frontmatter} describes each preliminary page and the rules that apply to it.

\section{Repository Layout}
\label{sec:layout}

\begin{description}
    \item[\file{thesis.tex}] The main file, and the only top-level \file{.tex} file you edit. It holds your class options, your thesis details, any extra packages, and the list of chapters.
    \item[\file{frontmatter/}] Your abstract, list of abbreviations, acknowledgements and dedication. All other preliminary pages are generated for you.
    \item[\file{content/}] One file per chapter and appendix. Images go in \file{content/figures/}. The \file{content/examples/} folder holds the table code shown in this document, and can be deleted.
    \item[\file{references.bib}] Your bibliography database.
    \item[\file{uvicphysastro.cls}] The document class. You should not need to edit it.
    \item[\file{extras/}] Bibliography style files (\file{.bst}), such as \file{aasjournal.bst} of the AAS journals (Chapter~\ref{ch:bibliographies}). Keep this folder next to \file{thesis.tex}.
\end{description}

\section{Quick Start}
\label{sec:quick-start}

\begin{enumerate}
    \item Download or clone the repository from \url{https://github.com/samdf96/UVic-Astronomy-LaTex-Template}, or upload all of its files to a new Overleaf project.
    \item In \file{thesis.tex}, set your class options (Section~\ref{sec:class-options}) and fill in your thesis details: \cmd{thesistitle}, \cmd{thesisauthor}, \cmd{thesisyear}, \cmd{thesisdegrees} and \cmd{thesiscommittee}.
    \item Write your abstract in \file{frontmatter/abstract.tex}, and edit (or delete) \file{frontmatter/abbreviations.tex}, \file{frontmatter/acknowledgements.tex} and \file{frontmatter/dedications.tex}.
    \item Replace the example chapters in \file{content/} with your own, and update the \cmd{include} lines in \file{thesis.tex} to match.
    \item Compile (Section~\ref{sec:compiling}).
\end{enumerate}

\section{Class Options}
\label{sec:class-options}

Class options are set in the \cmd{documentclass} line at the top of \file{thesis.tex}, for example:
\begin{lstlisting}[language=TeX]
\documentclass[degree=phd, bibliography=natbib]{uvicphysastro}
\end{lstlisting}
Any option you leave out takes its default value. Values are not case sensitive, and a misspelled value stops compilation with an error that lists the allowed values.

\begin{table}[htbp]
    \centering
    \caption{Class options of \texttt{uvicphysastro}. Defaults are listed first.}
    \label{tab:class-options}
    \begin{tabularx}{\linewidth}{@{}l l X@{}}
        \toprule
        Option & Values & Effect \\
        \midrule
        \texttt{degree} & \texttt{masters}, \texttt{phd} & ``Thesis'' and MASTER OF SCIENCE, or ``Dissertation'' and DOCTOR OF PHILOSOPHY, throughout the preliminary pages. \\
        \texttt{bibliography} & \texttt{none}, \texttt{natbib}, \texttt{biblatex} & Which bibliography package to load (Chapter~\ref{ch:bibliographies}). \\
        \texttt{bibstyle} & \texttt{authoryear}, \ldots & The \texttt{biblatex} citation style. Any \texttt{biblatex} style name is accepted. \\
        \texttt{spacing} & \texttt{onehalf}, \texttt{double} & Line spacing of the whole document. \\
        \texttt{floatfraction} & \texttt{0.5}, any number 0--1 & How full a page of floats must be before \LaTeX\ gives figures and tables a page of their own. Lower values (e.g.\ \texttt{0.01}) create float-only pages more readily. \\
        \texttt{draft} & \texttt{false}, \texttt{true} & Draft mode (Section~\ref{sec:draft-mode}). \\
        \texttt{listoftables} & \texttt{true}, \texttt{false} & Set to \texttt{false} if your thesis has no tables. \\
        \texttt{listoffigures} & \texttt{true}, \texttt{false} & Set to \texttt{false} if your thesis has no figures. \\
        \texttt{headings} & \texttt{left}, \texttt{centered} & Alignment of the preliminary-page headings (Section~\ref{sec:heading-styles}). \\
        \bottomrule
    \end{tabularx}
\end{table}

The page margins are 1~inch on all sides (with extra room at the top for the page number); UVic does not prescribe margins, so if you or your committee prefer a wider left margin, add a line such as \cmd{geometry}\texttt{\{left=1.5in\}} to the preamble of \file{thesis.tex}.

\section{Packages: What the Class Loads and What to Add}
\label{sec:packages}

The class loads every package that the thesis format itself depends on, together with a few general tools, so a thesis compiles without any \cmd{usepackage} lines of your own. Table~\ref{tab:class-packages} lists them. You do not need to load these again; if an older preamble of yours does, the repeated lines are harmless, as long as they do not ask for different options (see below).

\begin{table}[htbp]
    \centering
    \caption{Packages loaded by the class.}
    \label{tab:class-packages}
    \begin{tabularx}{\linewidth}{@{}>{\raggedright\arraybackslash}p{0.27\linewidth} >{\raggedright\arraybackslash}X@{}}
        \toprule
        Purpose & Packages \\
        \midrule
        Page layout & \texttt{geometry}, \texttt{setspace}, \texttt{fancyhdr}, \texttt{sectsty}, \texttt{tocloft} \\
        Mathematics & \texttt{amsmath}, \texttt{amsthm}, \texttt{amssymb} \\
        Figures and captions & \texttt{graphicx}, \texttt{caption}, \texttt{subcaption}, \texttt{adjustbox}, \texttt{pdfpages} \\
        Cross-references and links & \texttt{hyperref}, \texttt{varioref}, \texttt{url}, \texttt{xurl}, \texttt{notoccite} \\
        Bibliography & \texttt{natbib} or \texttt{biblatex}, depending on the \texttt{bibliography} option \\
        Symbols, colour and text & \texttt{wasysym}, \texttt{tipa}, \texttt{xcolor}, \texttt{xspace}, \texttt{verbatim} \\
        Draft mode only & \texttt{lineno}, \texttt{showkeys}, \texttt{changebar}, \texttt{datetime2}, \texttt{todonotes} \\
        \bottomrule
    \end{tabularx}
\end{table}

Anything else is up to you. Table~\ref{tab:optional-packages} lists packages that are often useful in physics and astronomy theses; none of them is required.

\begin{table}[htbp]
    \centering
    \caption{Optional packages that are often useful.}
    \label{tab:optional-packages}
    \begin{tabularx}{\linewidth}{@{}>{\raggedright\arraybackslash}p{0.38\linewidth} >{\raggedright\arraybackslash}X@{}}
        \toprule
        Package & Use \\
        \midrule
        \texttt{booktabs}, \texttt{multirow}, \texttt{xltabular}, \texttt{threeparttablex}, \texttt{pdflscape} & Tables (Chapter~\ref{ch:tables}); loaded in \file{thesis.tex} for the examples in this document \\
        \texttt{siunitx} & Numbers and units, e.g.\ \cmd{SI}\texttt{\{3.0e8\}\{m/s\}} \\
        \texttt{physics}, \texttt{slashed}, \texttt{esvect} & Physics notation: derivatives, bra--ket notation, Feynman slash, vector arrows \\
        \texttt{cleveref} & References that name what they point to, e.g.\ \cmd{cref}\texttt{\{fig:x\}} gives ``Figure 2.1'' \\
        \texttt{acronym}, \texttt{glossaries} & Lists of acronyms or symbols (Section~\ref{sec:abbreviations}) \\
        \texttt{listings} & Program code; used for the code examples in this document \\
        \texttt{fontspec} & Choosing fonts, with Lua\LaTeX\ only (Section~\ref{sec:engines}) \\
        \bottomrule
    \end{tabularx}
\end{table}

\subsection{Where to Load Packages}

Load your own packages in \file{thesis.tex}, in the block headed ``Your own packages and settings'', after the \cmd{documentclass} line. A few rules cover almost every case:
\begin{itemize}
    \item \textbf{Settings for packages the class loads go after \cmd{documentclass}}, in the same block. For example, \cmd{hypersetup}\texttt{\{hidelinks\}} makes all links black, and \cmd{captionsetup}\texttt{\{...\}} changes the caption style.
    \item \textbf{Options for packages the class loads go before \cmd{documentclass}.} Loading such a package again with different options stops compilation with an ``Option clash'' error. Instead, pass the options with \cmd{PassOptionsToPackage} on the very first line of \file{thesis.tex}; the class then loads the package with them. Numbered citations with \texttt{natbib} are the most common case (Section~\ref{sec:bib-numbered}):
\begin{lstlisting}[language=TeX]
\PassOptionsToPackage{numbers,sort&compress}{natbib}
\documentclass[bibliography=natbib]{uvicphysastro}
\end{lstlisting}
    \item \textbf{Your packages are loaded after \texttt{hyperref}}, because the class loads it. This is the order most packages expect, and some require it: \texttt{cleveref}, for example, must come after \texttt{hyperref}, which is automatically the case here.
\end{itemize}

\section{Compiling Your Thesis}
\label{sec:compiling}

The template was tested with \TeX\ Live 2026 (\LaTeX\ 2026-06-01; pdf\TeX\ 1.40.29, Lua\-HB\TeX\ 1.24.0, Xe\TeX\ 0.999998), using all three engines described in Section~\ref{sec:engines}. Older versions of \TeX\ Live should also work (most of the template was also tested with \TeX\ Live 2022), but only \TeX\ Live 2026 was tested in full.

\subsection{On Overleaf}

Upload every file in the repository to a new project, keeping the folder structure. In the project menu, set the main document to \file{thesis.tex} and the compiler to pdf\LaTeX\ (or Lua\LaTeX; see Section~\ref{sec:engines}). Overleaf runs the bibliography program and repeats the compilation as needed automatically. Also set the \TeX\ Live version, under \emph{Menu}~$\rightarrow$ \emph{Settings}~$\rightarrow$ \emph{\TeX\ Live version}, to 2026 or the newest version offered, so that your thesis is compiled with the version this template was tested with.

\subsection{On Your Own Computer}

You need a recent \TeX\ distribution: \TeX\ Live (Linux, Windows), Mac\TeX\ (macOS), or MiK\TeX\ (Windows). The simplest way to compile is \texttt{latexmk}, which runs \LaTeX\ and the bibliography program as many times as needed:
\begin{lstlisting}[language=bash]
latexmk -pdf thesis.tex
\end{lstlisting}
Most editors (VS Code with LaTeX Workshop, TeXstudio, TeXShop) use \texttt{latexmk} or an equivalent sequence for you. Without it, the full sequence for a \texttt{natbib} bibliography is \texttt{pdflatex}, \texttt{bibtex}, \texttt{pdflatex}, \texttt{pdflatex}. For \texttt{biblatex}, replace \texttt{bibtex} with \texttt{biber}. To use Lua\LaTeX, replace \texttt{pdflatex} with \texttt{lualatex}.

\subsection{Choosing a \TeX\ Engine}

\label{sec:engines}

The \TeX\ engine is the program that turns your \file{.tex} files into a PDF. Three engines are in common use, and the template has been tested with all of them: each produces the same thesis, page for page.

\begin{description}
    \item[pdf\LaTeX\ (recommended).] Use pdf\LaTeX\ unless you need one of the features listed under Lua\LaTeX\ below. It is the default on Overleaf, it compiles the fastest, and it is the engine that journal classes such as AASTeX and the MNRAS class are written for (and the only one arXiv accepts), so material moves between your papers and your thesis without surprises. Common accented characters such as \'e or \"o can be typed directly into your files; for anything more unusual, use the \LaTeX\ command for it.
    \item[Lua\LaTeX\ (fully supported).] Lua\LaTeX\ works with the template without any changes, so switch to it whenever you need to. It reads Unicode text natively, so you can type any character directly (non-Latin scripts, names with unusual diacritics, symbols), and it can use any font installed on your computer through the \texttt{fontspec} package. It is also where the newest \LaTeX\ developments arrive first, such as accessible (tagged) PDFs. The trade-off is speed: compilation takes several times longer than with pdf\LaTeX.
    \item[Xe\LaTeX.] Xe\LaTeX\ also works with the template, but it offers nothing that Lua\LaTeX\ does not, so choose Lua\LaTeX\ instead.
\end{description}

To change engine on Overleaf, open the project menu and change the \emph{Compiler} setting. To change it locally, give \texttt{latexmk} a different option:
\begin{lstlisting}[language=bash]
latexmk -pdf thesis.tex        # pdfLaTeX
latexmk -lualatex thesis.tex   # LuaLaTeX
\end{lstlisting}
Nothing in \file{thesis.tex} needs to change: the class detects the engine and sets up the fonts to suit it. Two things to keep in mind:
\begin{itemize}
    \item With Lua\LaTeX, do not load the \texttt{fontenc} or \texttt{inputenc} packages, which are common in older templates and examples; they are for pdf\LaTeX\ only. Choose fonts with the \texttt{fontspec} package instead, for example \cmd{setmainfont}\texttt{\{TeX Gyre Termes\}}.
    \item After switching engine, delete the auxiliary files (\texttt{latexmk -C}, or \emph{Recompile from scratch} on Overleaf) before compiling again.
\end{itemize}

\subsection{Compiling Selected Chapters}

A full thesis can take a while to compile. While you work on one chapter, uncomment the \cmd{includeonly} line in \file{thesis.tex} and list only the chapters you want:
\begin{lstlisting}[language=TeX]
\includeonly{content/chapter4_tables}
\end{lstlisting}
Page numbers, cross-references and citations to the other chapters are kept from the last full compilation. This only works for chapters added with \cmd{include} (not \cmd{input}).

\section{Draft Mode}
\label{sec:draft-mode}

The \texttt{draft=true} class option turns on several aids for writing and revising:
\begin{itemize}
    \item line numbers in the margin, so your committee can refer to exact lines;
    \item the key of every \cmd{label} and \cmd{ref}, printed in the margin;
    \item black bars marking lines that run into the margin (overfull boxes);
    \item a footer with the compile date and time on every page;
    \item margin notes with \cmd{todo}\texttt{\{...\}} (from the \texttt{todonotes} package);
    \item change bars: text between \cmd{cbstart} and \cmd{cbend} is marked in the margin, which is useful for showing revisions.
\end{itemize}
When \texttt{draft} is off, \cmd{todo} notes are silently hidden, so you can leave them in the source. Make sure draft mode is off for your final submission.

\section{Submitting Your Thesis}
\label{sec:submitting}

The final thesis is submitted as a PDF to UVicSpace. Before you submit, check the following:
\begin{itemize}
    \item \textbf{File name.} The PDF must be named \file{LastName_FirstName_Degree_Year.pdf}, for example \file{Smith_Jane_MSc_2026.pdf}, and the name must match the one on your title page. \LaTeX\ names the output after the main file (\file{thesis.pdf}), so rename the PDF before uploading it.
    \item \textbf{Draft mode is off} and no \cmd{todo} notes or change bars remain.
    \item \textbf{Page numbers in the table of contents are correct.} Compile twice after your last change; UVic reviewers spot-check them.
    \item \textbf{No signatures or personal contact information} appear anywhere, including the appendices.
    \item \textbf{The PDF is Unicode compliant.} The class takes care of this for pdf\LaTeX: all text in the PDF can be searched and copied.
\end{itemize}
Once a thesis is accepted to UVicSpace, it cannot be altered or resubmitted.
